\documentclass[10pt, letterpaper]{article}

% Packages:
\usepackage[
    ignoreheadfoot,
    top=1.2 cm,
    bottom=1.2 cm,
    left=1.6 cm,
    right=1.6 cm,
    footskip=1.0 cm,
]{geometry}
\usepackage{titlesec}
\usepackage{tabularx}
\usepackage{array}
\usepackage[dvipsnames]{xcolor}
\definecolor{primaryColor}{RGB}{0, 0, 0}
\definecolor{accentColor}{RGB}{40, 40, 40}
\usepackage{enumitem}
\usepackage{amsmath}
\usepackage[
    pdftitle={Yiyuan Li's CV},
    pdfauthor={Yiyuan Li},
    pdfcreator={LaTeX CV Yiyuan},
    colorlinks=true,
    urlcolor=accentColor
]{hyperref}
\usepackage[pscoord]{eso-pic}
\usepackage{calc}
\usepackage{bookmark}
\usepackage{lastpage}
\usepackage{changepage}
\usepackage{paracol}
\usepackage{ifthen}
\usepackage{needspace}
\usepackage{iftex}

\ifPDFTeX
    \input{glyphtounicode}
    \pdfgentounicode=1
    \usepackage[T1]{fontenc}
    \usepackage[utf8]{inputenc}
    \usepackage{charter}
\else
    % XeTeX/tectonic: charter.sty has no TU fonts, so load XCharter (OpenType) for real bold/italic
    \usepackage{fontspec}
    \setmainfont{XCharter}[
        Extension=.otf,
        UprightFont=*-Roman,
        BoldFont=*-Bold,
        ItalicFont=*-Italic,
        BoldItalicFont=*-BoldItalic
    ]
\fi

\newcommand{\githublink}[2]{\href{#1}{#2 $\nearrow$}}

% Some settings:
\raggedright
\AtBeginEnvironment{adjustwidth}{\partopsep0pt}
\pagestyle{empty}
\setcounter{secnumdepth}{0}
\setlength{\parindent}{0pt}
\setlength{\topskip}{0pt}
\setlength{\columnsep}{0.15cm}
\pagenumbering{gobble}

\titleformat{\section}{\needspace{4\baselineskip}\bfseries\large}{}{0pt}{}[\vspace{1pt}\titlerule]

\titlespacing{\section}{-1pt}{0.3 cm}{0.15 cm}

\renewcommand\labelitemi{$\vcenter{\hbox{\small$\bullet$}}$}
\newenvironment{highlights}{
    \begin{itemize}[
        topsep=0.08 cm,
        parsep=0.05 cm,
        partopsep=0pt,
        itemsep=0pt,
        leftmargin=0.4 cm
    ]
}{
    \end{itemize}
}

\newenvironment{onecolentry}{
    \begin{adjustwidth}{0.15 cm}{0.15 cm}
}{
    \end{adjustwidth}
}

\newenvironment{twocolentry}[2][]{
    \onecolentry
    \def\secondColumn{#2}
    \setcolumnwidth{\fill, 4.5 cm}
    \begin{paracol}{2}
}{
    \switchcolumn \raggedleft \secondColumn
    \end{paracol}
    \endonecolentry
}

\newenvironment{header}{
    \setlength{\topsep}{0pt}\par\kern\topsep\centering\linespread{1.5}
}{
    \par\kern\topsep
}

\let\hrefWithoutArrow\href

\begin{document}
    \newcommand{\AND}{\unskip
        \cleaders\copy\ANDbox\hskip\wd\ANDbox
        \ignorespaces
    }
    \newsavebox\ANDbox
    \sbox\ANDbox{$|$}

    \begin{header}
        \fontsize{25 pt}{25 pt}\selectfont Yiyuan Li

        \vspace{5 pt}

        \normalsize
        \mbox{Singapore}%
        \kern 5.0 pt%
        \AND%
        \kern 5.0 pt%
        \mbox{\hrefWithoutArrow{mailto:yy@eyuan.me}{yy@eyuan.me}}%
        \kern 5.0 pt%
        \AND%
        \kern 5.0 pt%
        \mbox{\hrefWithoutArrow{tel:+6589523447}{+65 8952 3447}}%
        \kern 5.0 pt%
        \AND%
        \kern 5.0 pt%
        \mbox{\hrefWithoutArrow{https://www.eyuan.me/}{eyuan.me}}%
        \kern 5.0 pt%
        \AND%
        \kern 5.0 pt%
        \mbox{\hrefWithoutArrow{https://linkedin.com/in/yyuanl}{linkedin.com/in/yyuanl}}%
        \kern 5.0 pt%
        \AND%
        \kern 5.0 pt%
        \mbox{\hrefWithoutArrow{https://github.com/yuann3}{github.com/yuann3}}%
    \end{header}

    \vspace{5 pt - 0.3 cm}

    \section{Summary}
        \begin{onecolentry}
            Software engineer building AI agent infrastructure: agent tools and the sandboxing around them, evals, and the quality layer under inference routing that keeps multi-tenant agent platforms reliable. I anchor my work in the fundamentals of software engineering, which is why my side projects rebuild the classics: a SQLite engine from scratch, a Redis server, a GPU shader renderer that runs inside Emacs buffers, and an actor-based multi-agent runtime in Rust. 2+ years of experience, mostly in Rust, C, and TypeScript, with Python and Go when the job calls for it. I'm happiest figuring out how something works one layer deeper than I need to. I think computers and Emacs are cool, and so are house music and movies.
        \end{onecolentry}

    \section{Technical Skills}
        \begin{onecolentry}
            \textbf{Languages:} Rust, C, C++, TypeScript, Python, Go (Golang), SQL, Emacs Lisp, GLSL, Bash
        \end{onecolentry}

        \vspace{0.07 cm}

        \begin{onecolentry}
            \textbf{AI Agents \& Inference:} MCP, OAuth 2.1, Tool Design \& Approval Gates, LLM-as-Judge \& Probe Evals, Inference Provider Routing, Streaming \& Cancellation, Mastra AI, RAG
        \end{onecolentry}

        \vspace{0.07 cm}

        \begin{onecolentry}
            \textbf{Frameworks \& Libraries:} Tokio, wgpu, Bun, Hono, React, Vue 3, FastAPI, Flask, Node.js, Playwright
        \end{onecolentry}

        \vspace{0.07 cm}

        \begin{onecolentry}
            \textbf{Databases \& Infra:} PostgreSQL (RLS, LISTEN/NOTIFY), pg-boss, Redis, SQLite, ChromaDB, Supabase, AWS, Docker, Turborepo, Git, Linux, CI/CD
        \end{onecolentry}

        \vspace{0.07 cm}

        \begin{onecolentry}
            \textbf{Concepts:} Agent Runtimes, Multi-Tenancy, Actor Model, Distributed Systems, Concurrency, Database Internals, GPU Programming, FFI/Dynamic Modules, TDD
        \end{onecolentry}

    \section{Experience}

        \begin{twocolentry}{
            Jun 2026 -- Present
        }
            \textbf{Software Engineer}, Voltade Pte. Ltd. -- Singapore
        \end{twocolentry}
        \vspace{0.08 cm}
        \begin{onecolentry}
            \begin{highlights}
                \item Built the inference-quality layer under Volty's multi-provider LLM routing: detectors score every production reply for broken output, hosts that underperform their peers are quarantined, and a streaming breaker retries looping replies on another host
                \item Built the eval tooling around it: a probe harness that tests inference endpoints on tool calls, structured JSON, and long context against a reference endpoint, plus the LLM-judge eval behind smart inbox filters
                \item Designed and built agent tools: the first version of org MCP servers (OAuth 2.1, per-tool approval, hash-pinned definitions), PDF operations, and web fetch behind an SSRF guard with DNS pinning
                \item Split agent turns into foreground and background concurrency tiers, streamed tokens to the browser over Postgres NOTIFY and SSE with tenant scoping that fails closed, and made Stop a hard abort of in-flight tool calls
                \item Tech Stack: TypeScript, Bun, Hono, PostgreSQL (RLS), pg-boss, React 19, MCP, Playwright, AWS
            \end{highlights}
        \end{onecolentry}

        \vspace{0.15 cm}

        \begin{twocolentry}{
            Jan 2026 -- May 2026
        }
            \textbf{Software Engineer -- AI}, ClassDo Pte. Ltd. -- Singapore
        \end{twocolentry}
        \vspace{0.08 cm}
        \begin{onecolentry}
            \begin{highlights}
                \item Designed and built a multi-agent AI curriculum platform from scratch with a small team, and owned the architecture for how agents coordinate, validate, and hand off work
                \item Built the backend in a Turborepo monorepo: Hono middleware with JWT auth, RBAC, a typed RPC client, and 809+ tests across 59 files, deployed on AWS
                \item Integrated Singapore's SkillsFuture Skills Framework into Supabase (PostgreSQL), seeding 269K rows across 10 tables with multi-country schema support
                \item Built the Vue 3 frontend with Vike SSR, Pinia, and PrimeVue, with Storybook for the design system
                \item Tech Stack: TypeScript, Mastra AI, Hono, Vue 3, Supabase, PostgreSQL, AWS, Turborepo, Vitest, Zod
            \end{highlights}
        \end{onecolentry}

        \vspace{0.15 cm}

        \begin{twocolentry}{
            Jan 2025 -- Sep 2025
        }
            \textbf{Software Engineering Intern}, Newcastle Australia IHE Pte. Ltd. -- Singapore
        \end{twocolentry}
        \vspace{0.08 cm}
        \begin{onecolentry}
            \begin{highlights}
                \item Built a full-stack RAG learning platform with FastAPI, React 19, and Ollama, giving 500+ users question answering over their course documents
                \item Built the document pipeline on LlamaParse with ChromaDB vector search, cutting query response time by 40\%
                \item Built OAuth 2.0 authentication and REST API endpoints handling 1,000+ daily requests
                \item Tech Stack: Python, TypeScript, React 19, Next.js, FastAPI, Flask, Tailwind CSS, JWT, ChromaDB
            \end{highlights}
        \end{onecolentry}

        \vspace{0.15 cm}

        \begin{twocolentry}{
            Feb 2024 -- Nov 2024
        }
            \textbf{Technical Leader, Data Structures \& Algorithms}, University of Newcastle
        \end{twocolentry}
        \vspace{0.10 cm}
        \begin{onecolentry}
            \begin{highlights}
                \item Led algorithm workshops on dynamic programming state transitions, graph traversal edge cases, and backtracking with pruning
                \item Wrote teaching material for the hard parts: B-tree rebalancing, amortized complexity, and off-by-one errors in binary search
            \end{highlights}
        \end{onecolentry}

    \section{Projects}

        \begin{twocolentry}{
            \githublink{https://github.com/yuann3/mra}{\textbf{View on GitHub}}
        }
            \textbf{MRA: Multi-Agent Runtime for Rust}
        \end{twocolentry}
        \vspace{0.08 cm}
        \begin{onecolentry}
            \begin{highlights}
                \item Designed an Erlang/OTP-inspired agent runtime where each AI agent runs as a Tokio actor with bounded mpsc mailboxes, zero shared mutable state, and transparent restarts via ArcSwap mailbox slots
                \item Implemented a supervisor with OneForOne/OneForAll restart strategies, exponential backoff, hang detection, and peer injection -- agents survive crashes without callers noticing
                \item Technologies: Rust, Tokio, async/await, Actor Model, Serde, Figment, OpenRouter API
            \end{highlights}
        \end{onecolentry}

        \vspace{0.15 cm}

        \begin{twocolentry}{
            \githublink{https://github.com/yuann3/emaclaude}{\textbf{View on GitHub}}
        }
            \textbf{Emaclaude: Claude Code Orchestrator}
        \end{twocolentry}
        \vspace{0.08 cm}
        \begin{onecolentry}
            \begin{highlights}
                \item Built a 3-agent coding workflow (planner, coder, reviewer) inside Doom Emacs split buffers, coordinated by a \textasciitilde130-line Rust CLI and an Elisp orchestrator piping JSON state and effects
                \item Implemented a review loop with double-confirmation passes to catch hallucinated approvals, plus inline diff comments and one-command PR creation via Magit and gh CLI
                \item Technologies: Rust, Emacs Lisp, Claude Code CLI, Magit, Serde, Clap, GitHub CLI
            \end{highlights}
        \end{onecolentry}

        \vspace{0.15 cm}

        \begin{twocolentry}{
            \githublink{https://github.com/yuann3/shaderview}{\textbf{View on GitHub}}
        }
            \textbf{ShaderView: GLSL Shader Engine for Emacs}
        \end{twocolentry}
        \vspace{0.08 cm}
        \begin{onecolentry}
            \begin{highlights}
                \item Built a GPU-accelerated GLSL shader renderer that runs inside Emacs buffers using wgpu (Vulkan/Metal), with hot-reload on save and Shadertoy-compatible uniforms (iResolution, iTime, iMouse)
                \item Wrote a Rust dynamic module (cdylib FFI) bridging Emacs Lisp to a pure-Rust GPU core, supporting concurrent shader sessions, PNG export, and configurable frame rates
                \item Technologies: Rust, wgpu, Emacs Lisp, GLSL, FFI/Dynamic Modules, Vulkan, Metal
            \end{highlights}
        \end{onecolentry}

        \vspace{0.15 cm}
    
        \begin{twocolentry}{
            \githublink{https://github.com/yuann3/RayRust}{\textbf{View on GitHub}}
        }
            \textbf{RayRust: CPU/GPU Path Tracer}
        \end{twocolentry}
        \vspace{0.08 cm}
        \begin{onecolentry}
            \begin{highlights}
                \item Built a Monte Carlo path tracer in Rust with Lambertian, metal, and dielectric (glass) materials, and a camera with depth of field and adjustable field of view
                \item Added a GPU backend on wgpu with WGSL shaders behind a feature flag, so the same scenes render on Vulkan, Metal, or DirectX 12, with the CPU renderer as fallback
                \item Technologies: Rust, wgpu, WGSL, GPU Compute, Monte Carlo Sampling
            \end{highlights}
        \end{onecolentry}

        \vspace{0.15 cm}

        \begin{twocolentry}{
            \githublink{https://github.com/yuann3/sequel/}{\textbf{View on GitHub}}
        }
            \textbf{Sequel: SQLite Database Engine}
        \end{twocolentry}
        \vspace{0.08 cm}
        \begin{onecolentry}
            \begin{highlights}
                \item Built a SQLite-compatible database engine from scratch in Rust, implementing B-tree data structures and binary file parsing to read .db files without external dependencies
                \item Implemented SQL query execution including SELECT with WHERE clauses, COUNT aggregation, and index optimization achieving 3x faster lookups through B-tree traversal
                \item Technologies: Rust, B-tree, Binary Parsing, Database Internals, File I/O
            \end{highlights}
        \end{onecolentry}

        \vspace{0.15 cm}
    
        \begin{twocolentry}{
            \githublink{https://github.com/yuann3/Rego/}{\textbf{View on GitHub}}
        }
            \textbf{Rego: Redis Server Implementation}
        \end{twocolentry}
        \vspace{0.08 cm}
        \begin{onecolentry}
            \begin{highlights}
                \item Implemented Redis-compatible server in Go with core commands (GET, SET, EXPIRE) and WAIT replication, passing 100\% of Codecrafters integration tests
                \item Optimized memory allocation using slice pooling, reducing garbage collection overhead and improving throughput by 30\%
                \item Technologies: Go (Golang), TCP/IP Networking, Concurrency, Memory Optimization
            \end{highlights}
        \end{onecolentry}

        \vspace{0.15 cm}
    
        \begin{twocolentry}{
            \githublink{https://github.com/yuann3/Ruskey/}{\textbf{View on GitHub}}
        }
            \textbf{Ruskey: Programming Language Interpreter}
        \end{twocolentry}
        \vspace{0.08 cm}
        \begin{onecolentry}
            \begin{highlights}
                \item Developed a complete interpreter for Monkey language in Rust: lexer, recursive descent parser, AST, and tree-walking evaluator
                \item Implemented closures, first-class functions, and hash maps using Test-Driven Development (TDD) methodology
                \item Technologies: Rust, Compiler Design, Abstract Syntax Trees (AST), TDD
            \end{highlights}
        \end{onecolentry}

        \vspace{0.15 cm}

    \section{Education}

        \begin{twocolentry}{
            Nov 2025 -- Sep 2026
        }
            \textbf{Singapore University of Technology and Design (SUTD)}, 42 Singapore Computer Science
        \end{twocolentry}  
        \vspace{0.08 cm}
        \begin{onecolentry}
            \begin{highlights}
                \item Core Curriculum: Peer-to-peer software engineering program focused on C, Unix, and systems programming
                \item Completed 16+ projects: libft, ft\_printf, get\_next\_line, push\_swap, pipex, Born2beRoot
            \end{highlights}
        \end{onecolentry}

        \vspace{0.15 cm}

        \begin{twocolentry}{
            Jan 2024 -- Sep 2025
        }
            \textbf{The University of Newcastle}, Bachelor of Information Technology
        \end{twocolentry}  
        \vspace{0.08 cm}
        \begin{onecolentry}
            \begin{highlights}
                \item High Distinction: Object-Oriented Programming (OOP)
                \item Distinction: Data Structures and Algorithms (DSA), Advanced Database Systems
            \end{highlights}
        \end{onecolentry}

\end{document}
